\documentclass[11pt,a4paper]{article}
\usepackage[utf8]{inputenc}
\usepackage[T1]{fontenc}
\usepackage{amsmath,amssymb,mathtools}
\usepackage{geometry}
\usepackage{hyperref}
\geometry{margin=2.5cm}
\title{QCE, Measurement Independence, Delayed Choice and the Formal Boundary of Superdeterminism}
\author{Ingolf Lohmann}
\date{6 August 2026}

\begin{document}
\maketitle

\begin{abstract}
This note separates three claims that are often conflated in discussions of delayed-choice experiments and superdeterminism: (i) later measurement choice does not imply a physical rewriting of an already registered past; (ii) measurement independence excludes the measurement-dependent escape route usually called superdeterminism in Bell reasoning; and (iii) spacelike separation and local response structure do not by themselves establish measurement independence.  The QIK-VRT Lean formalization contains a finite common-cause countermodel for (iii) and a conditional kernel theorem for (ii).  Consequently, a physically relevant QCE exclusion of superdeterminism still requires an independently justified derivation and empirical/reference binding of measurement independence.
\end{abstract}

\section{Epistemic claim boundary}
Let $\lambda$ denote a hidden preparation state and $a,b$ the two measurement settings.  The standard Bell-style measurement-independence condition is
\[
P(\lambda\mid a,b)=P(\lambda).
\]
At support level this is represented by factorization of jointly admissible hidden-state/setting triples:
\[
J(\lambda,a,b)\iff H(\lambda)\land S(a,b).
\]
The QIK-VRT kernel calls this predicate \texttt{MeasurementIndependent}.

A measurement-dependent Bell escape route is represented by
\[
\neg\operatorname{MeasurementIndependent}(M).
\]
The formal proxy \texttt{SuperdeterministicCandidate} is intentionally restricted to that dimension; it is not claimed to encode every philosophical use of ``superdeterminism''.

\section{Conditional exclusion theorem}
By definition,
\[
\operatorname{MeasurementIndependent}(M)
\Longrightarrow
\neg\operatorname{SuperdeterministicCandidate}(M).
\]
The Lean theorem is
\begin{verbatim}
measurementIndependence_excludes_superdeterministicCandidate
\end{verbatim}
and has the type
\[
\forall M,\; MI(M)\to\neg SD_{MI}(M).
\]
This theorem is logically exact but conditional.  It does not establish $MI$ for nature.

\section{Why delayed choice and spacelike separation are insufficient}
The formal two-wing model gives each wing only its local setting:
\[
A=A(\lambda,a),\qquad B=B(\lambda,b).
\]
Therefore changing $b$ cannot directly change the response function on wing A, and changing $a$ cannot directly change the response function on wing B.  This captures the relevant structural locality at the response stage.

Nevertheless define a finite common-cause model with a hidden bit $\lambda\in\{0,1\}$ and admissibility
\[
J(\lambda,a,b)\iff(a=\lambda)\land(b=\lambda),
\]
while every hidden value and every setting pair are individually admissible.  Then local response structure holds, but support factorization fails.  In Lean this is the theorem
\begin{verbatim}
localResponseStructure_not_sufficient_for_measurementIndependence
\end{verbatim}
with the accompanying theorem
\begin{verbatim}
commonCauseModel_not_measurementIndependent.
\end{verbatim}
Thus spacelike separation of the two setting events rules out direct signal exchange between the wings at choice time, but not a common past cause correlating $\lambda,a,b$.

\section{Delayed-choice interpretation}
The original delayed-choice quantum-eraser literature reports correlations sorted according to a later measurement basis.  Standard quantum mechanics does not require the later basis choice to rewrite a previously registered local event.  Operationally, the later comparison changes which conditional correlation is evaluated, not the historical detector record itself.

The QIK-VRT epistemic reading is therefore
\[
\text{later verification}\neq\text{past mutation},
\]
while
\[
\text{later verification}=\text{new evidence constraining justified classification}.
\]

\section{The QCE proof obligation}
The substantial QCE claim would have to establish the first implication in
\[
\boxed{
QCE\text{-axioms}
\Longrightarrow MI
\Longrightarrow \neg SD_{MI}
}.
\]
The second implication is kernel-formalized.  The first is not licensed merely by delayed choice, locality of the response signatures, or the experimenters' assertion of free choice.

The Lean structure \texttt{QCEFreedomCertificate} therefore records the outstanding obligation.  Given such a certificate, Lean proves
\begin{verbatim}
qceFreedomCertificate_excludes_superdeterministicCandidate.
\end{verbatim}
Promoting the certificate from a formal assumption to a physical claim requires an interpretation/reference map and evidence, in accordance with the existing QIK-VRT world-formula relation kernel.

\section{Scientific interpretation}
The result establishes two kernel-level facts:
\begin{enumerate}
\item measurement independence is sufficient to exclude the formal measurement-dependent superdeterministic candidate;
\item local two-wing response structure is not sufficient to derive measurement independence.
\end{enumerate}
It does not establish that nature satisfies measurement independence, and therefore does not by itself establish that nature is non-superdeterministic.

This is consistent with the standard Bell-theorem literature, which explicitly identifies measurement independence as a supplementary assumption and superdeterminism as one route that rejects it.

\section{Sources}
\begin{itemize}
\item Stanford Encyclopedia of Philosophy, ``Bell's Theorem'', \url{https://plato.stanford.edu/entries/bell-theorem/}.
\item Y.-H. Kim et al., ``Delayed `Choice' Quantum Eraser'', Phys. Rev. Lett. 84, 1 (2000), DOI: \url{https://doi.org/10.1103/PhysRevLett.84.1}.
\item R. Chaves, G. B. Lemos, J. Pienaar, ``Causal Modeling the Delayed-Choice Experiment'', Phys. Rev. Lett. 120, 190401 (2018), DOI: \url{https://doi.org/10.1103/PhysRevLett.120.190401}.
\item The BIG Bell Test Collaboration, ``Challenging local realism with human choices'', Nature 557, 212--216 (2018), DOI: \url{https://doi.org/10.1038/s41586-018-0085-3}.
\item S. Hossenfelder and T. N. Palmer, ``Rethinking Superdeterminism'', arXiv:1912.06462.
\item QIK-VRT pre-spacetime ontology publication receipt, DOI: \url{https://doi.org/10.5281/zenodo.21804399}.
\end{itemize}

\section{Repository bindings}
The formal module is
\begin{verbatim}
formalization/QIKVRT_Formalization_v2.0/
QIKVRTFormalization/QuantumFoundations/MeasurementIndependence.lean
\end{verbatim}
with a dedicated \texttt{AxiomAudit.lean}.  A successful Lean 4.19 execution-bound receipt is required before any theorem is described as kernel-verified for the final publication candidate.

\section{Conclusion}
The strongest currently justified statement is
\[
\boxed{MI\Rightarrow\neg SD_{MI}}
\]
combined with the finite countermodel
\[
\boxed{\text{local response structure}\not\Rightarrow MI}.
\]
Therefore a non-circular QCE proof against superdeterminism must derive measurement independence from deeper QCE premises rather than assume it under the name of free choice.

\bigskip
\noindent\textit{Quod erat demonstrandum---for the formal logical boundary.}

\noindent Ingolf Lohmann
\end{document}
